% Abhakseite „Das kann ich“ – Zone nur im Gesamt (\mitzone)
%% TODO Ich-kann-Sätze: Platzhalter wie in den Aufgabendateien
\begin{abhakseite}
\ifdefined\mitzone
\abhakgruppe{Kennst du schon}
\abhak{1}{Quadrieren auch negativer Zahlen und Dezimalzahlen, Quadrat vor Punkt vor Strich: (−7)² = 49, nicht −49; 0,5² = 0,25}
\abhak{2}{Koordinaten lesen und eintragen, vier Quadranten, (x | y)-Reihenfolge, Kästchenraster mit 0,5-Einteilung}
\abhak{3}{Lineare Funktion f(x) = m·x + n: Gerade zeichnen, Funktionswert, Punktprobe, Nullstelle}
\abhak{4}{Binomische Formel ausmultiplizieren und zusammenfassen, auch mit Minus in der Klammer}
\abhak{5}{Lineare Gleichung lösen, Terme mit x auf eine Seite bringen, Schreibform mit Strich}
\abhak{6}{Schnittpunkt zweier Geraden durch Gleichsetzen}
\abhak{7}{Quadratwurzel ziehen, auch mit dem Taschenrechner, Näherungswert runden; aus einer negativen Zahl gibt es keine Wurzel}
\abhak{8}{Quadratische Gleichung durch Wurzelziehen lösen und Normalform mit der p-q-Formel lösen, beide Lösungen angeben}
\abhak{9}{Binomische Formel ausmultiplizieren und zusammenfassen, auch mit Minus in der Klammer – Fehler finden}
\abhak{10}{Binomische Formel ausmultiplizieren und zusammenfassen, auch mit Minus in der Klammer – selbst rechnen}
\fi
\abhakgruppe{Einheit 1 · Normalparabel und Streckfaktor}
\abhak{11}{Normalparabel und Streckfaktor}
\abhak{12}{Normalparabel und Streckfaktor – Fehler finden, Begründen, Darstellungswechsel, Anwendung}
\abhakgruppe{Einheit 2 · Scheitelpunktform}
\abhak{13}{Scheitelpunktform}
\abhak{14}{Scheitelpunktform – Fehler finden, Begründen, Darstellungswechsel, Anwendung}
\abhakgruppe{Einheit 3 · Normalform}
\abhak{15}{Normalform}
\abhak{16}{Normalform – Fehler finden, Begründen}
\abhakgruppe{Einheit 4 · Nullstellen und Schnittpunkte berechnen}
\abhak{17}{Nullstellen und Schnittpunkte}
\abhak{18}{Nullstellen und Schnittpunkte – Fehler finden, Begründen, Darstellungswechsel, Anwendung}
\end{abhakseite}
